\documentclass[10pt,a4paper]{amsart}
\usepackage[margin=19mm]{geometry}
\usepackage{amsmath,amssymb,mathtools}
\usepackage[scr=rsfs,cal=euler]{mathalfa}
\usepackage{xfrac}
\usepackage[unicode,hidelinks,bookmarks=false]{hyperref}
\newcommand{\ie}{i.e.\ }
\newcommand{\cf}{cf.\ }
\newcommand{\resp}{respectively\ }
\newcommand{\Subsection}{\subsection}
\providecommand{\nobreakdash}{\nobreak\hbox{-}}
\setlength{\emergencystretch}{2em}
\multlinegap0pt
\usepackage{bm}
\usepackage{bbm}
\usepackage{dsfont}
\usepackage{xspace}
\usepackage{cancel}
\usepackage{mathtools}
\usepackage{amsthm}
\makeatletter
\@ifundefined{theorem}{\newtheorem{theorem}{Theorem}}{}
\@ifundefined{corollary}{\newtheorem{corollary}[theorem]{Corollary}}{}
\@ifundefined{lemma}{\newtheorem{lemma}[theorem]{Lemma}}{}
\@ifundefined{prop}{\newtheorem{prop}[theorem]{Proposition}}{}
\makeatother
\newcommand{\Int}{\mathrm{Int\,}}
\newcommand{\Op}{ {\mathcal O}{\it p}\,}
\newcommand{\R}{\mathbb{R}}
\newcommand{\Skel}{\mathrm{Skel}}
\newcommand{\eps}{\epsilon}
\newcommand{\p}{\partial}
\newcommand{\wh}{\widehat}
\title[Honda-Huang's work on contact convexity revisited]{Honda-Huang's work on contact convexity revisited}
\author{Yakov Eliashberg, Dishant Pancholi}
\date{Source: arXiv:2207.07185v2 (CC BY 4.0)}
\begin{document}
\maketitle

\noindent\emph{Source and licence.} Honda-Huang's work on contact convexity revisited. Version arXiv:2207.07185v2 (CC BY 4.0), distributed under the Creative Commons Attribution 4.0 International licence, as recorded in the arXiv licence metadata for this exact version and shown on the abstract page https://arxiv.org/abs/2207.07185v2. Full rights notice: licenses/arxiv-2207.07185-CC-BY-4.0.txt.\par\smallskip

\noindent\textit{Editorial scope.} Counted unit: the Prop whose source label is \texttt{prop:prelim-plug} in the article (source0.tex lines 1032--1052, source label \texttt{prop:prelim-plug}), delivered together with the article's own complete original proof of it (source0.tex lines 1112--1151, one \texttt{proof} environment of 40 lines). No mathematics is written, repaired or re-derived: statement and proof are the article's own text line for line. The proof is detached from the statement in the source: the article states the result at lines 1032--1052 and prints its proof at lines 1112--1151, and the proof environment's own header names the statement, so the association is the article's own. Required context retained uncounted: lm:Lyapunov-criterion (lines 183--191); cor:good-Lyapunov (lines 402--403); lm:2d-plug (lines 873--917); lm:widehatX (lines 1070--1086). Every definition, display and label of those lines is retained unchanged. Omitted from this package and not redistributed: 1--182, 192--401, 404--872, 918--1031, 1053--1069, 1087--1111, 1152--1624. Nothing inside a retained excerpt is omitted or altered: every retained body is delivered exactly as the article's lines are written, so no omission receipt of any kind accompanies this package. Citations: the retained excerpts carry no citation command at all, so no bibliography entry of the article is transcribed and no reference is redistributed. Reference adaptation: none was needed and none was made; every reference inside the delivered unit resolves inside it, so no unresolved reference or citation is produced. Printed slips kept literal: no printed word, symbol, hypothesis or number of the counted statement or of its proof was altered. Layout and class: the article's own class is \texttt{amsart} with packages including amssymb, amsmath, amsthm, amsbsy, amscd, bm, graphicx, xcolor, verbatim; the wrapper is the lane's pinned \texttt{amsart} editorial wrapper at 10pt on A4, which restates the theorem-like environments the retained text needs and adds the article's own macro definitions that the retained text needs (\texttt{\textbackslash Int}, \texttt{\textbackslash Op}, \texttt{\textbackslash R}, \texttt{\textbackslash Skel}, \texttt{\textbackslash eps}, \texttt{\textbackslash p}, \texttt{\textbackslash wh}), with extra wrapper packages (bm, bbm, dsfont, xspace, cancel, mathtools). The delivered numbering is the unit's own. Assets: no retained span contains a figure environment, a graphics-inclusion command, a PDF-page inclusion, a table or a TikZ picture; no image file is copied into this package, the package declares no asset dependency and no image file is redistributed with it. Licence: version arXiv:2207.07185v2 of this article is distributed under the Creative Commons Attribution 4.0 International licence, as recorded in the arXiv licence metadata for this exact version and shown on the abstract page https://arxiv.org/abs/2207.07185v2. A full rights notice is delivered at licenses/arxiv-2207.07185-CC-BY-4.0.txt. Characters: every retained excerpt is pure ASCII, so the delivered engine, XeLaTeX with its default Latin Modern text font, reports no missing character.
\begin{lemma}\label{lm:Lyapunov-criterion}

 Let $X$ be a vector field  with isolated hyperbolic zeroes which are non-degenerate or of embryo type on a closed $m$-dimensional manifold $\Sigma$.
Then $X$ admits a Lyapunov function if and only if the following conditions are satisfied:
\begin{itemize}
\item[(L1)]   every trajectory of $X$ originates and terminates at a  zero of $X$;
\item[(L2)]  there exists an  ordering  $O_1,\dots, O_N$ of zeroes such that there are no trajectories of $X$ which originate at $O_i$ and terminate in $O_j$ if $i>j$.
\end{itemize}
\end{lemma}

 \begin{corollary}\label{cor:good-Lyapunov}
 Suppose that a vector field $X$ satisfies conditions (L1) and (L2) from   Lemma \ref{lm:Lyapunov-criterion}.  Then it admits a good Lyapunov function if and only if it  has no retrograde connections. In particular, any  $X$ which satisfies (L1) and the   Morse-Smale condition admits a good Lyapunov function. \end{corollary}

\begin{lemma}\label{lm:2d-plug} For any $  \eps>0$
there exists an      isotopy 
$h_s:R\to O$, $s\in[0,1] $, which is fixed together with its $\infty$-jet along $\p R$,   constant for $s\in[0,\frac18]$,  and such that the following properties  a)--i) are satisfied.
Denote $\beta_s:=h_s^* (dz + x dy)$. Let $Y_s$ be the vector field directing the characteristic foliation $\ell_s$ of $\beta_s$.  
\begin{itemize}
   \item[a)] $d\beta$ restricted to the interior $ \Int R_{+}$ of $R_ +$  is positive, and  $ d\beta$ restricted to the interior $\Int R_{-}$ of $R_-$ is negative
for all   $s\in(\frac18,1]$;
\item[b)] $\beta_s$ has
\begin{itemize}
\item no zeros    for $s<\frac34$,
\item  a positive and negative embryos $o_+:=(\frac16,\psi_1(\frac16)),
o_-=(\frac56,\psi_2(\frac56))$ for $s=\frac34$,
 \item a pair 
 \begin{align*}
 &e_+(s)=(e_+^1(s),\psi_1(e_+^1(s))), \hbar_+(s)=(\hbar_+^1(s),\psi_1(\hbar_+^1(s))),\; 0<e_+^1(s)<\hbar_+^1(s)<\frac13
 \end{align*} of positive elliptic and  hyperbolic   points, and a pair
 \begin{align*}
  &e_-(s)=(e_-^1(s),\psi_1(e_-^1(s))), \hbar_-(s)=(\hbar_-^1(s),\psi_1(\hbar_-^1(s))),\; 1>e_+^1(s)>\hbar_+^1(s)>\frac23
 \end{align*} of negative elliptic and  hyperbolic points for $s>\frac34$.
\end{itemize}
 
\item[c)] the incoming separatrices of $\hbar_+(s)$   and $o_+$  for $Y_s$,    are  contained in $\Psi_1$, and outgoing separatrices
of $\hbar_-(s)$   and $o_-$  for $Y_s$ are  contained in $\Psi_2$,  $s\geq 
\frac34$;
 \item[d)] there exists $\eps_1\in(0,\eps)$ such that the outgoing separatrices of $\hbar_+(s)$ for $Y_1$ terminate at $e_-$ and $(1,\eps_1 )$, and the   incoming separatrices of $\hbar_-(s)$ for $Y_1$ originate at $e_+$ and $(0,1-\eps_1)$; 
\item[f)]   $Y_s$  for $s\in[\frac12,1]$ is outward transverse  to the graph 
$\Theta$, viewed as a  part of the boundary of the domain $R_+$;
\item[g)]   $Y_{s}$ admits   a  family of good Lyapunov function  $\psi_{s}:R\to\R$   such that  $\psi_{s}|_{\Op\p R}=y$;
\item[h)] $\beta_s(\frac{\p}{\p z})>0$  everywhere in $R$ for $s\in[\frac12,1]$;
\item[i)]  for any $\sigma>0$  the isotopy $h_s$, $s\in[0,1]$,  can be extended to a  2-parametric isotopy $h_{s,t}, 0\leq s,t\leq1$,  such that
\begin{itemize}
\item $h_{s,0}=h_s$, $h_{0,t}=h_0$ for all $s,t\in[0,1]$;
\item $h_{s,t}=h_{s}$   for $s\leq\frac34-\sigma, t\in[0,1]$;
\item for each $s\in(\frac34,1]$  the isotopy $h_{s,t},t\in[0,1]$, is supported in a    $\sigma$-neighborhood of the separatrices  connecting $\hbar_\pm(s)$ with $e_\pm(s)$;   for each $s\in[\frac34-\sigma,1]$ $h_{s,t}$ is supported in a    $\sigma$-neighborhood of $o_\pm$;  
 \item $h_{s,t}$ is   $\sigma$-close in the $C^0$-sense to $h_{s,0}$ for all $s,t\in[0,1]$;
\item  the family of vector fields  $Y_{s,t}$ directing the characteristic foliations $\ell_{s,t}$ of $\beta_{s,t}:=h_{s,t}^*\alpha$ admits a  family  good Lyapunov functions  $\psi_{s,t}:R\to\R$   such that  $h_{s,t}|_{\Op\p R}=y$;
\item $Y_{s,t}$ has a pair  of positive elliptic and hyperbolic zeroes at
$e_\pm(s(1-2t) +\frac{3t}2)$ and $\hbar_\pm(s(1-2t) +\frac{3t}2)$ for $s>\frac34, t<\frac12$, pairs of embryos at $o_\pm$ for $s=\frac34, t=\frac12$ and no zeroes otherwise;
\item $h_{s,t}=h_{0}$ for $s\in[0,\frac18], t\in[0,1],$ and $h_{s,t}=h_{1,t}$ for $s\in[\frac78,1]$    $t\in[0,1]$;
\item $X_{s,t}$ is outwardly transverse to $\Theta$ for all $s\in[\frac12,1],\; t\in[0,1]$.
  \end{itemize}
 
\end{itemize}

\end{lemma}

 \begin{lemma}\label{lm:widehatX}
 We have 
\begin{itemize}
\item[1.] $X= Y+aZ$ over $W_0\times Q$, where the function $a:Q\to\R$ is determined by the equation 
\begin{align}\label{eq:function-a1}
a\mu=\iota(Y)d\mu.
\end{align}
\item[2.] $X=Y+aZ+bR$ over $C\times Q$, where the functions
  $a,b:Q\to\R$ are determined  by the equations
  \begin{equation}\label{lm:widehatX2}
  \begin{split}
 & a\kappa_\tau=\iota(Y_\tau)d\mu_\tau, \\
 & b\tau+\mu_\tau(Y_\tau)=0.
 \end{split}
 \end{equation}
 \end{itemize}
 \end{lemma}

\begin{prop}\label{prop:prelim-plug}
 The vector field   $X_{s,t}$ on $Q$, $s,t \in[0,1] $, has the following properties.  
\begin{enumerate}
 
\item  every trajectory of $X_{1}$
\begin{enumerate}
\item  which starts  at
 ${}^{in}P  $ converges to a zero of $X_{1}$;

 \item  which starts  at a point $(u,z,0) \in {}^{in}T$    exits   at a point $(u',z',1)\in {}^{out}S   $,   where  $u'=Z^{-a}(u) , a>0$  and  $z'\in(1-\frac\eps2,1]$;  
\item   which ends 
 ${}^{out}P $ originates at a zero of $X_{1}$;
 \item which terminates  at a point $(u,z,1) \in {}^{out}T $ begins at   a point $(u',z',0)  \in {}^{in}S \subset {}^{in}V, $  where  $u'=Z^{-a}(u) , a>0$, and  $z'\in[0,\frac\eps2)$;  
\end{enumerate}
 
\item  vector fields $X_{s,1}$, $s\in[0,1] $, have no zeroes;
\item $\wh\Gamma_+$ is invariant with respect to the negative flow of $X_1$, and $\wh\Gamma_-$ is invariant with respect to the positive one;
\item all trajectories of $X_s$ which converge to positive singularities either do not intersect $\p Q$ and are contained in $\wh\Gamma_+$, or intersect it at points of $\Skel(W)\times\frac\eps2\subset {}^{in}V$; all trajectories of $X_s$ which originate at negative singularities either do not intersect $\p Q$ and are contained in $\wh \Gamma_-$, or intersect it at points of $\Skel(W)\times(1-\frac{\eps}2)\subset {}^{out}V;$
\item  the  family of vector fields $X_{s,t}$ admits a family of good Lyapunov functions $ \psi_{s,t}$, equal to $y$ on $\p Q$;
\end{enumerate}
\end{prop}  

\begin{proof}[Proof of Proposition \ref{prop:prelim-plug}]
(1a) When a   trajectory  of $X_1$ which enters at a point of  $ {}^{in}P$ it is in  the region where  $X_1=Y_1 + a Z$ with $a<0$. Since $a$ is negative the trajectory continues to remain in the region where the plug is given by $ W_0\times R,$  and hence projects onto trajectories of $-Z$ and $Y_1$, when projected to the corresponding factors. Moreover, the projection to $R$ remains in $R_-$, according to Property f) of Lemma \ref{lm:2d-plug}, 
 and therefore remains  in the region where the coefficient $a$ is negative.
 But in $R_-$ any trajectory of $Y_1$ entering at a point in $[\eps,1-\eps_1)\times 0$ converges to a negative  zero of $Y_1$, while every trajectory of $-Z$ converges to a zero of $Z$. Hence, any trajectory of $Y_1$ 
 entering through $ {}^{in}P$ converges to a $0$ of the vector field $X_1$.
 \medskip
 
 (1b)
  If a trajectory enters  at a  $(u, z)\in{}^{in}T  $ then similarly to 1a) we have  $X_1=Y_1+aZ$ for a negative coefficient $a$, and therefore, it projects    onto trajectories of $-Z$ and $Y_1$ in the factors $W$ and $R$ respectively. But every trajectory of $Y_1$ which enters through $(1-\eps_1,1]\times 0$ exits $R$ at a point of $[1-\frac\eps2,1]$, and hence the corresponding trajectory of $X_1$ exits ${}^{out}V$ at a point $(u'=Z^{-c}(u), z')$ for some positive $c>0$ and $z'\in(-\frac\eps2.1]$.
  \medskip
  
 (1c) and (1d) follows from the same arguments as, respectively, (1a) and (1b) applied to the vector field $-X_1$.
 
 \medskip
 Properties (2)-(4) are straightforward from the corresponding properties in Lemmas \ref{lm:2d-plug} and  \ref{lm:widehatX}.
 
(5) According to  Corollary \ref{cor:good-Lyapunov} it is sufficient to verify that
\begin{itemize}
\item each trajectory of $X_{s,t}$  either originates at $V\times 0$, or at a critical point of $X_{s,t}$;
\item  each trajectories of $X_{s,t}$  either terminates at $V\times 1$, or at a critical point of $X_{s,t}$;
\item $X_{s,t}$ has no retrograde connections.
\end{itemize}
In addition to the Weinstein subdomain $W_0=W\setminus\p W\times(1-\eps,1]$ consider also a larger subdomain  $W_1=W\setminus\p W\times(1-\frac\eps2,1]$.
Denote $Q_\pm:=W_1\times R_\pm,  Q_b:=(W\setminus W_1)\times R\subset Q.$  We have $Q=Q_+\cup Q_-\cup Q_b$.
Let us first analyze  the forward trajectory $X_{s,t}^u(p)=(w(u),r(u))$, $u\in\R_+$ of a  point 
 $p=(w,r)\in Q$. 
 
 If $p\in Q_-$ and $w\in W_0$ then $w(u)$  belongs to the negative Liouville trajectory $\bigcup\limits_{\tau\geq 0}Z^{-\tau}(w)$ as long as $r(u)\in R_-$.
 Similarly, if $w\in W_1\setminus W_0=\p W\times[1-\eps,1-\frac\eps2]$ the second coordinate of $w$ decreases as long as $r(u)\in R_-$.
But  the $R$-component $Y_{s,t}$  of $X_{s,t}$ is by construction inwardly transverse to the boundary of $R_-$, and hence,  remains in $R_-$ for all $u\geq 0$. Therefore, the trajectory either converges to a singular point of $X_{s,t}$, or exits through $V\times 1$.

Suppose $p\in Q_b$. Recall that the  $R$-component $Y_{s,t}$  of $X_{s,t}$ has in $Q_b$ a positive projection to  the  $y$-direction. Hence,  $X_{s,t}^u(p)$  either  exits through $V\times 1$, or enters $Q_+\cup Q_-$. But  in that case it only can enter $Q_-$, and therefore, the analysis of the previous case does apply.
 
 Suppose now that $p\in Q_+$.    If $w\in W_0$ then $w(u)$ moves along a positive Liouville trajectory of $w$, and if  $w\in \p W\times [1-\eps ,1-\frac\eps2]$,  then the second coordinate of $w(u)$ increases   as long as $r(u)\in R_+$. Hence, the trajectory either exits through $V\times 1$, or  enters either $Q_b$ or $Q_-$, and therefore, the previous analysis applies.
 
  Backward trajectories could be analyzed  similarly, with exchanging $Q_+$ and $Q_-$ cases.
 The absence of retrograde connections follows from (4).
 

 \end{proof}

\end{document}
